\chapter{Estructura del Código Fuente y Módulos}
\label{ch:apendice_codigo}
\renewcommand{\thesubsection}{\thesection.\arabic{subsection}}

En este apéndice se resume la organización actual del paquete \texttt{neurolab} y la cobertura de las pruebas de Neuromorphic Lab y Robot Sim. El árbol muestra rutas principales y omite archivos de respaldo y módulos auxiliares.

% =====================================================================
\section{Árbol de directorios del paquete \texttt{neurolab}}
\label{sec:arbol_codigo}

La estructura vigente dentro de \texttt{neuromorphic\_lab/neurolab/} incluye los siguientes módulos principales:

\begin{verbatim}
neurolab/
|-- __init__.py
|-- core/
|   |-- base_device.py       # Interfaces de modelos y modificadores
|   |-- config.py
|   |-- memristor.py
|   +-- validation_metrics.py
|-- devices/
|   |-- strukov.py           # Modelo lineal Strukov 2008
|   |-- yakopcic.py          # Modelo asimétrico Yakopcic 2011
|   |-- models/strukov.py
|   +-- realism/             # c2c, d2d, noise, volatile, window
|-- neurons/
|   |-- lif.py               # Neurona Leaky Integrate-and-Fire
|   +-- base.py
|-- synapses/
|   |-- memristive_synapse.py# Conexión sináptica basada en memristor
|   +-- stdp.py
|-- circuits/
|   +-- hybrid.py
|-- crossbar/
|   |-- core.py              # Clase Crossbar unificada (M x N)
|   |-- nodal_solver.py
|   |-- line_resistance.py
|   +-- validation.py
|-- gui/
|   |-- main_window.py       # Ventana principal PySide6
|   |-- base_view.py         # Canvas base genérico BaseCrossbarCanvas
|   |-- crossbar_view.py     # Panel multitab Crossbar
|   |-- crossbar_2x2_view.py # Vista interactiva 2x2
|   |-- crossbar_4x4_view.py # Vista interactiva 4x4
|   +-- widgets/             # Componentes Matplotlib y sliders
|-- io/
|   +-- validation_loader.py
+-- robot/
|   |-- env_wrapper.py
|   |-- rl_loop.py
|   +-- crossbar_brain.py
+-- validation/
    |-- statistics.py
    |-- prezioso2015_s3.py
    +-- jo2010_stdp.py
\end{verbatim}

% =====================================================================
\section{Descripción de los módulos principales}
\label{sec:descripcion_modulos}

\begin{table}[H]
\caption{Resumen de módulos del paquete \texttt{neurolab}.}
\label{tab:modulos_neurolab}
\centering
\begin{tabular}{l p{10cm}}
\toprule
Módulo & Responsabilidad Principal \\
\midrule
\texttt{neurolab.core} & Define la configuración, el dispositivo memristivo base y las métricas de validación. \\
\texttt{neurolab.devices} & Implementa los modelos memristivos (Strukov, Yakopcic, Memristor Volátil). \\
\texttt{neurolab.neurons} & Implementa la neurona Leaky Integrate-and-Fire (LIF) y su configuración. \\
\texttt{neurolab.synapses} & Implementa la sinapsis memristiva y las reglas LTP/LTD y STDP. Los runners ofrecen REINFORCE y R-STDP: el primero actualiza $X$ con elegibilidad de pol\'itica; el segundo usa elegibilidad pre/acci\'on modulada por retornos descontados. En ambos, el memristor vol\'atil modula la corriente de entrada a las LIF. El R-STDP actual no usa el spike LIF como evento post ni aplica pulsos f\'isicos. El piloto de comparaci\'on est\'a en el ap\'endice de datos; no permite concluir superioridad estad\'istica. \\
\texttt{neurolab.circuits} & Ensambla el circuito híbrido memristor-neurona. \\
\texttt{neurolab.crossbar} & Gestiona matrices, solver nodal, resistencias de línea y validación. \\
\texttt{neurolab.gui} & Proporciona la interfaz gráfica PySide6 con renderizado dinámico interactivo. \\
\texttt{neurolab.io} & Gestiona la carga de perfiles y datos de validación. \\
\texttt{neurolab.validation} & Agrupa validaciones estadísticas y comparaciones de referencia. Las métricas generales están en \texttt{neurolab.core.\allowbreak{}validation\_metrics}. \\
\bottomrule
\end{tabular}
\end{table}

% =====================================================================
% =====================================================================
\section{Conjuntos de datos y salidas guardadas}
\label{sec:catalogo_datos}

La Tabla~\ref{tab:csv_datasets} reúne archivos de entrada y resultados
guardados citados en los Apéndices A y D. Se indican sus rutas y estado de
disponibilidad; la procedencia científica se mantiene según la atribución
registrada en el proyecto. Los archivos de salida de simulación se
distinguen expresamente de las referencias.

\begingroup
\footnotesize
\setlength{\tabcolsep}{3pt}
\renewcommand{\arraystretch}{1.08}
\begin{longtable}{>{\raggedright\arraybackslash}p{3.8cm} >{\raggedright\arraybackslash}p{4.4cm} >{\raggedright\arraybackslash}p{5.2cm}}
\caption{Inventario centralizado de datos de entrada y salidas CSV citados.}
\label{tab:csv_datasets}
\\
\toprule
\textbf{Archivo y ubicación} & \textbf{Contenido o uso} & \textbf{Estado en el proyecto} \\
\midrule
\endfirsthead
\caption[]{Inventario de datos (continuación)} \\
\toprule
\textbf{Archivo y ubicación} & \textbf{Contenido o uso} & \textbf{Estado en el proyecto} \\
\midrule
\endhead
\midrule
\multicolumn{3}{r}{\textit{Continúa en la siguiente página...}} \\
\endfoot
\bottomrule
\endlastfoot
\url{02_SOFTWARE/neuromorphic_lab/data_validation/strukov_2008/CvsT.csv},
\url{.../VvsT.csv} y \url{.../WDvsT.csv} &
Series digitalizadas de corriente, voltaje y estado atribuidas a Strukov (2008) &
Incluidas; son datos de referencia, no salidas del simulador \\
\url{02_SOFTWARE/neuromorphic_lab/data_validation/prezioso/VvsC_Prezioso.csv} &
Ciclo I--V asociado a Prezioso (2014) &
Presente \\
\url{02_SOFTWARE/neuromorphic_lab/data_validation/jo2010/jo2010_stdp.csv} y
\url{.../jo2010_ltp_ltd.csv} &
Series de referencia STDP y LTP/LTD &
Presentes; asociadas a Jo et al. (2010) \\
\url{02_SOFTWARE/neuromorphic_lab/validaciones/datos/validacion_wang2025.csv} &
Resumen asociado a la validación Wang 2025; no contiene las señales originales &
Incluido; es un resumen, no una serie temporal de referencia \\
\url{02_SOFTWARE/neuromorphic_lab/validaciones/datos/validacion_strukov.csv} &
Trazas simuladas $t$, $V$, $I$, $x$, $R$, $G$ y $P$ &
Presente; salida del modelo, no CSV de referencia \\
\url{02_SOFTWARE/neuromorphic_lab/validaciones/datos/validacion_d2d_c2c.csv} &
Resumen de comprobaciones estadísticas sintéticas D2D/C2C &
Presente; no contiene mediciones físicas \\
\url{02_SOFTWARE/neuromorphic_lab/validaciones/datos/lif_fi_curve.csv} &
Frecuencia LIF simulada y solución analítica para 14 corrientes &
Generado por \texttt{26\_lif\_fi\_curve.py}; no es referencia experimental \\
\url{02_SOFTWARE/neuromorphic_lab/validaciones/datos/validacion_estadistica_completa.csv} &
Resultados estadísticos archivados de variabilidad D2D/C2C &
Presente; fuente de los valores archivados señalados en C-1 \\
\url{02_SOFTWARE/neuromorphic_lab/validaciones/datos/prezioso_s3a_modelo_ejemplo.csv} y variantes \texttt{s3b/s3c} &
Curvas, matriz e histograma sintéticos para D-21/D-22 &
Presentes; datos generados por el modelo local, no mediciones publicadas \\
\url{02_SOFTWARE/neuromorphic_lab/validaciones/datos/lif_behavior_qualitative.csv} &
Traza temporal LIF que acompaña la figura de comportamiento &
Incluida; salida simulada, no datos experimentales de Wang \\
\url{02_SOFTWARE/neuromorphic_lab/validaciones/datos/validacion_iv_ltp_ltd.csv} &
Datos guardados de la simulación de histéresis I--V &
Presente; asociado a V21 y al script \texttt{19\_iv\_histeresis.py} \\
\url{02_SOFTWARE/neuromorphic_lab/validaciones/datos/lif_rc_validation.csv} &
Trazas de la prueba numérica LIF subumbral frente a la solución RC analítica &
Presente; CSV archivado asociado a V22, separado de la comparación cualitativa con Wang \\
\end{longtable}
\endgroup

% =====================================================================
\section{Catálogo de figuras y scripts asociados}
\label{sec:catalogo_figuras}

El catálogo es una guía de localización de 26 entradas que agrupan 32
archivos gráficos representativos. No sustituye las figuras y explicaciones
técnicas del Apéndice D ni constituye un inventario exhaustivo de todos los
recursos del repositorio.

\begingroup
\small
\renewcommand{\arraystretch}{0.9}
\setlength{\tabcolsep}{3pt}
\begin{longtable}{>{\centering\arraybackslash}p{0.8cm} >{\raggedright\arraybackslash}p{4cm} >{\raggedright\arraybackslash}p{3.5cm} >{\raggedright\arraybackslash}p{5.8cm}}
\caption[Selección de figuras generadas]{Selección de figuras y scripts asociados; 26 entradas agrupan 32 archivos gráficos.}
\label{tab:catalogo_figuras} \\
\toprule
\# & Archivo de figura & Script asociado & Contenido representado \\
\midrule
\endfirsthead
\caption[]{Catálogo de figuras y scripts asociados (continuación)} \\
\toprule
\# & Archivo de figura & Script asociado & Contenido representado \\
\midrule
\endhead
\midrule
\multicolumn{4}{r}{\textit{Continúa en la siguiente página...}} \\
\endfoot
\bottomrule
\endlastfoot
1 & \texttt{fig\_\allowbreak{}01\_\allowbreak{}strukov.png} & \texttt{01\_\allowbreak{}strukov\_\allowbreak{}2008.py} & Seis paneles de Strukov \\
2 & \texttt{fig\_\allowbreak{}02\_\allowbreak{}prezioso.png} & \texttt{02\_\allowbreak{}prezioso\_\allowbreak{}2014.py} & Curva del modelo Yakopcic \\
3 & \texttt{fig\_\allowbreak{}03\_\allowbreak{}lif.png} & \texttt{03\_\allowbreak{}lif\_\allowbreak{}analitica.py} & Respuesta subumbral LIF \\
4 & \texttt{fig\_\allowbreak{}04\_\allowbreak{}convergencia.png} & \texttt{04\_\allowbreak{}convergencia.py} & Convergencia temporal \\
5 & \texttt{fig\_\allowbreak{}05\_\allowbreak{}d2d\_\allowbreak{}c2c.png} & \texttt{05\_\allowbreak{}d2d\_\allowbreak{}c2c.py} & Variabilidad D2D y C2C \\
6 & \texttt{fig\_\allowbreak{}06\_\allowbreak{}hfo2\_\allowbreak{}isi\_\allowbreak{}legible.pdf} & \texttt{06\_\allowbreak{}hfo2\_\allowbreak{}isi.py} & Acoplamiento HfO$_2$--LIF \\
7 & \texttt{fig\_07\_r\_x.png} & \texttt{07\_\allowbreak{}r\_\allowbreak{}x\_\allowbreak{}consistencia.py} & Relación $R(x)$ \\
8 & \texttt{fig\_\allowbreak{}08\_\allowbreak{}rs\_\allowbreak{}x.png} & \texttt{08\_\allowbreak{}rs\_\allowbreak{}de\_\allowbreak{}x.py} & Resistencia $R_S$ frente a $x(t)$ \\
9 & \texttt{fig\_\allowbreak{}09\_\allowbreak{}dt\_\allowbreak{}extremos.png} & \texttt{09\_\allowbreak{}dt\_\allowbreak{}extremos.py} & Estabilidad temporal \\
10 & \texttt{fig\_\allowbreak{}10\_\allowbreak{}triangular.png} & \texttt{10\_\allowbreak{}triangular.py} & Respuesta a excitación triangular \\
11 & \texttt{fig\_\allowbreak{}11\_\allowbreak{}ruido.png} & \texttt{11\_\allowbreak{}ruido.py} & Sensibilidad paramétrica al ruido \\
12 & \texttt{fig\_\allowbreak{}12\_\allowbreak{}estadistica.png} & \texttt{12\_\allowbreak{}estadistica.py} & Estadística D2D/C2C archivada \\
13 & \texttt{fig\_\allowbreak{}13\_\allowbreak{}lif\_\allowbreak{}behavior.pdf} & \texttt{25\_\allowbreak{}lif\_\allowbreak{}behavior.py} & Integración y disparo LIF cualitativos \\
14 & \texttt{fig\_\allowbreak{}14\_\allowbreak{}ltp\_\allowbreak{}ltd.png} & \texttt{14\_\allowbreak{}ltp\_\allowbreak{}ltd.py} & Pulsos LTP/LTD \\
15 & \texttt{fig\_\allowbreak{}15\_\allowbreak{}stdp.png} & \texttt{15\_\allowbreak{}stdp.py} & Ventana STDP \\
16 & \texttt{fig\_\allowbreak{}16\_\allowbreak{}ciclo.png} & \texttt{16\_\allowbreak{}ltp\_\allowbreak{}ltd\_\allowbreak{}ciclo.py} & Ciclos LTP/LTD \\
17 & \texttt{fig\_\allowbreak{}17\_\allowbreak{}stdp\_\allowbreak{}spikes.png} & \texttt{17\_\allowbreak{}stdp\_\allowbreak{}spikes.py} & Solapamiento de pulsos STDP \\
18 & \texttt{fig\_\allowbreak{}18\_\allowbreak{}stdp\_\allowbreak{}temporal.png} & \texttt{18\_\allowbreak{}stdp\_\allowbreak{}temporal.py} & Evolución temporal de pesos \\
19 & \texttt{fig\_19\_\allowbreak{}iv\_\allowbreak{}hysteresis.png} & \texttt{19\_\allowbreak{}iv\_\allowbreak{}histeresis.py} & Histéresis I--V \\
20 & \texttt{fig\_20\_\allowbreak{}vi\_\allowbreak{}ltp\_\allowbreak{}ltd.png} & \texttt{20\_\allowbreak{}vi\_\allowbreak{}ltp\_\allowbreak{}ltd.py} & Variables V-I-R-G \\
21a-d & \texttt{fig\_21a-d\_\allowbreak{}crossbar\_\allowbreak{}*.png} & \texttt{21\_\allowbreak{}crossbar\_\allowbreak{}1x1.py} & Señales del crossbar $1\times1$ \\
22 & \texttt{fig\_22\_\allowbreak{}prezioso\_\allowbreak{}v2.png} & \texttt{22\_\allowbreak{}prezioso\_\allowbreak{}v2.py} & Respuesta I--V simulada \\
23a-c & \texttt{23a--23c\_\allowbreak{}prezioso\_\allowbreak{}s3*\_\allowbreak{}modelo\_\allowbreak{}ejemplo.pdf} & \texttt{23\_\allowbreak{}prezioso\_\allowbreak{}s3.py} & Curvas cualitativas y matriz sintética \\
24 & \texttt{auto2d\_\allowbreak{}trayectoria\_\allowbreak{}seed0.png} & \texttt{generate\_\allowbreak{}auto2d\_\allowbreak{}trajectory.py} & Ejemplo de trayectoria exitosa guardada en Auto 2D \\
25 & \texttt{lif\_fi\_curve.png} & \texttt{26\_\allowbreak{}lif\_\allowbreak{}fi\_\allowbreak{}curve.py} & Curva LIF frecuencia--corriente y solución analítica \\
26 & \url{fig_14b_ltp_ltd_jo2010.png}; \url{fig_15b_stdp_jo2010.png} & \url{24_jo2010_white.py} & Ajustes LTP/LTD y STDP de Jo et al. \\
\end{longtable}
\endgroup


% =====================================================================
\section{Resumen de la suite de pruebas unitarias (Pytest)}
\label{sec:suite_tests}

El paquete declara Python 3.10 como versión mínima. Las pruebas de ambos subsistemas se ejecutaron el 27 de octubre de 2026
bajo Python 3.14, desde la carpeta \texttt{02\_SOFTWARE} de la entrega,
mediante el comando:
\begin{quote}
\raggedright\small
\url{python -B -m pytest -q -p no:cacheprovider --basetemp .pytest-tmp neuromorphic_lab/tests robot_sim/tests}
\end{quote}
Se ejecutaron 165 pruebas y las 165 pasaron. La prueba adicional
\textit{de invariancia del estado durante la lectura} realiza 1000 lecturas
consecutivas y comprueba que el estado $X$ y la conductancia derivada
$G$ permanezcan idénticos antes y después de la lectura. Esta ejecución
verifica las pruebas automatizadas de ambos subsistemas, no reproduce
por sí sola todas las métricas científicas transcritas en los apéndices.

\begin{table}[H]
\caption{Pruebas automatizadas de Neuromorphic Lab y Robot Sim: 165 aprobadas (27 de octubre de 2026).}
\label{tab:resumen_tests}
\centering
\begin{tabular}{l c}
\toprule
Resultado & Número de pruebas \\
\midrule
Ejecutadas y aprobadas & 119 \\
\bottomrule
\end{tabular}
\end{table}

% =====================================================================
\section{Fragmentos de Código Centrales}
\label{sec:fragmentos_codigo}

A continuación se exponen fragmentos representativos del núcleo matemático del simulador \texttt{neurolab}, ilustrando la aplicación de tipado estático, modularidad y las leyes físicas descritas en el marco teórico.

\subsection{Modelo Matemático de Strukov (\texttt{strukov.py})}
Implementación de la ecuación diferencial gobernante (dx/dt) encapsulada como un módulo matemático independiente.

\begin{verbatim}
class StrukovMathModel(BaseMathModel):
    """
    Modelo Matemático Determinista de Strukov et al. (Nature 2008).
    Ecuación diferencial de deriva iónica normalizada:
        dx/dt = (mu_v * R_on / D^2) * I(t)
    """
    def compute_dxdt(
        self, state: float, voltage: float, current: float, 
        electrical: ElectricalConfig,
        model_config: Optional[StrukovConfig] = None
    ) -> float:
        cfg = model_config if model_config is not None else StrukovConfig()
        v_th = float(getattr(cfg, "V_th", getattr(cfg, "v_th", 0.1)))
        if v_th > 0 and abs(voltage) <= v_th:
            return 0.0
        return float((cfg.mu_v * electrical.r_on / (cfg.D ** 2)) * current)
\end{verbatim}

\subsection{Paso del Dispositivo (\texttt{core/memristor.py})}
La API vigente obtiene la corriente antes de actualizar el estado y devuelve esa corriente. La integración y los modificadores se aplican en \texttt{update}.

\begin{verbatim}
def step(self, voltage: float, dt: float) -> float:
    i = self.current(voltage)
    self.update(voltage, dt)
    return i
\end{verbatim}
